\chapter{Human authority and the Chrome extension}
\label{chap:governance-interface}

\section{Purpose and scope}

The enforcement mechanisms in \Cref{chap:design} need an operator interface
because some safe outcomes cannot be represented as an automatic allow or
block. A policy may permit an action only after narrow consent; a protected
value may need human entry without authorising the model to type it; or the
observation boundary may be unable to establish that the next model input is
clean. In each case, the interface is part of the authority protocol rather
than a notification added after the decision.

\systemname\ implements that protocol as a Manifest V3 Chrome extension backed
by the local Sentinel service. The extension has four operational surfaces:

\begin{enumerate}
  \item a \textbf{side panel} containing the task conversation, live run state,
  intervention cards, and pause, resume, stop, and steering controls;
  \item an \textbf{in-page control layer} that makes agent ownership visible,
  offers cooperative input shielding, and supports direct takeover;
  \item an \textbf{Activity dashboard} that presents completed runs,
  privacy-preserving enforcement records, and recorded
  intervention-timing breakdowns;
  and
  \item \textbf{model and policy settings} that configure the local runtime.
\end{enumerate}

A separate one-time server-hosted page supports manual unpacked installation.
The implemented thesis scope comprises approval, takeover,
unresolved-observation incident, trace, and extension UI primitives with
server-owned lifecycle contracts. Together they form the implemented
authority protocol and interface system described here. The backend lifecycle
contracts are exercised by tests; the client is type-checked and packaged by
the production build.
Operator-visible plan authority, provenance hiding, and a multi-agent
supervisory queue are separate future mechanisms, not features attributed to
this implementation. Broader human-factors evaluation remains incomplete: no
participant study or WASP oracle result evaluates operator performance.

\section{Extension architecture}

\subsection{Separation between presentation and authority}

The Python service remains the security authority. It owns the run, policy,
pending requests, timers, audit sequence, and privileged executor. The
extension renders state and submits bounded responses; client state cannot
authorise an action because the blocked coroutine waits on a server request.

\Cref{fig:extension-architecture} shows the principal control paths. The side
panel and dashboard receive an initial snapshot followed by sequenced events
over a localhost WebSocket. User commands return through a localhost REST API.
The extension service worker maintains a compact control projection in
\texttt{chrome.storage.session}; the content script reads that projection to
render controls inside the active page. This avoids giving the page access to
the full conversation or trace. The content script uses a closed Shadow DOM so
page styles do not accidentally alter its internals. The host still sits in the
untrusted document and can be hidden, removed, or imitated by page code; the
overlay is therefore advisory presentation, not a tamper-resistant boundary.

\begin{figure}[htbp]
\centering
\begin{tikzpicture}[
  node distance=9mm and 13mm,
  box/.style={draw,rounded corners,align=center,minimum height=10mm,
    minimum width=31mm,fill=white},
  trusted/.style={box,draw=sentinelblue,very thick},
  ui/.style={box,draw=sentinelamber},
  page/.style={box,draw=sentinelred},
  flow/.style={-{Latex[length=2mm]},thick},
  both/.style={<->,>=Latex,thick},
  weak/.style={-{Latex[length=2mm]},dashed}
]
  \node[trusted] (server) {Sentinel service\\policy and authority};
  \node[ui,right=32mm of server] (panel) {Side panel and\\dashboard};
  \node[ui,below=of server] (worker) {Extension\\service worker};
  \node[ui,below=of worker] (content) {Closed-shadow\\page controls};
  \node[page,below=of content] (web) {Active web page\\and human input};

  \draw[both] (server) --
    node[above,font=\scriptsize]{snapshot + events}
    node[below,font=\scriptsize]{bounded responses} (panel);
  \draw[both] (server) -- node[right,align=left,font=\scriptsize]
    {compact state\\REST commands} (worker);
  \draw[flow] (worker) -- node[right,font=\scriptsize]{session state} (content);
  \draw[weak] (content) -- node[right,font=\scriptsize]{shield and controls} (web);
\end{tikzpicture}
\caption{Chrome extension control architecture. Security decisions and pending
authority remain in the local service; extension components present state and
submit constrained responses. The active page is untrusted; the in-page
extension component is presentation and transport, not authority.}
\label{fig:extension-architecture}
\end{figure}

The service binds to \texttt{127.0.0.1}. Its origin middleware accepts
Chrome-extension origins and the service's own localhost origins for ordinary
API requests, validates JSON fields and request identifiers, and permits the
in-page handover endpoint to resolve only an existing opaque handover
identifier. These checks prevent an arbitrary website from using the normal
control API, but localhost binding is not an authentication boundary against
another compromised process running as the same operating-system user. That
process and a compromised browser or extension are outside the accepted thesis
boundary.

\subsection{Packaging and browser lifecycle}

The Vite build packages the React and TypeScript client and its fonts as a
static Manifest V3 bundle. The browser manager launches a dedicated persistent
Chrome profile and reconnects through its saved DevTools endpoint. It probes
for a hello beacon or extension target, then tries launch-flag and DevTools
loading. If branded Chrome rejects both, a one-time localhost page supplies the
exact \emph{Load unpacked} path and confirms the next beacon. Another Chrome
may install the bundle as a remote control and monitor; privileged execution
remains in the dedicated governed browser.

\subsection{Snapshot and event consistency}

Two failures are especially dangerous in a human-authority interface:
displaying a stale request after it has expired, and losing a live request
during a panel reload. \systemname\ addresses both with server-owned pending
state. A WebSocket client first receives a snapshot containing run status,
recent events, configuration, pending approvals, and pending handovers. It
then applies monotonic sequenced events. The side panel retains a bounded
200-event live feed and reconciles the active conversation with persisted
turns.

On disconnection, the client enters an explicit offline state and reconnects
with exponential backoff from approximately 1 to 8 seconds plus jitter. The
next snapshot replaces the pending-card lists, so a resolved request
disappears and an unresolved one reappears. The extension service worker uses
the same snapshot-first rule for the in-page projection. If several control
requests are present, an unresolved security incident has display priority so
that an ordinary handover cannot conceal it.

This protocol separates visual from authority state. Client expiry hides a
stale card, but the server timer owns the timeout. A response only resolves the
request identifier; execution still requires any fresh preflight and its
terminal audit event.

Sequence numbers are process-monotonic, but replay and per-client queues are
bounded and the client does not diagnose sequence gaps. A slow client can
therefore miss non-authority telemetry. Pending approvals and handovers are
reconstructed separately from server-owned managers, so a telemetry gap cannot
mint consent or erase a live authority request.

\section{Authority taxonomy}

\systemname\ uses four distinct outcomes because ``ask the user'' is too vague
to be a security primitive. The division allocates decision selection and
action implementation separately rather than choosing one blanket level of
automation~\cite{parasuraman2000automation}.
\Cref{tab:authority-taxonomy} specifies what each path means and, critically,
what it does not authorise.

\begin{table}[htbp]
\centering
\caption{Authority outcomes exposed by the Chrome extension.}
\label{tab:authority-taxonomy}
\small
\begin{tabularx}{\textwidth}{@{}L{0.17\textwidth}L{0.23\textwidth}Y@{}}
\toprule
Outcome & Operator choice & Security meaning \\
\midrule
Hard refusal &
No approval control &
The proposed effect is prohibited. The agent action does not execute and a
human response cannot override the rule through the approval interface. \\

Action approval &
Allow once, bounded similar scope, or deny &
The named pending proposal may proceed only after server validation and a
fresh preflight. Consent is not a general instruction to trust the page. \\

Direct takeover &
Perform the step in the page, optionally add a note, then return control &
The human, not the agent, performs a protected or intrinsically human step.
Returning control does not authorise an intercepted agent disclosure or
guarantee that later page observations omit the entered value. \\

Security incident &
Review or change the page, request a fresh observation, or stop &
The unresolved observation is discarded. Re-checking does not approve hidden
content and cannot resume from the quarantined input. \\
\bottomrule
\end{tabularx}
\end{table}

Ordinary supervision is separate from these enforcement outcomes. While a run
is working, the operator may steer the agent at its next reasoning boundary,
pause and take manual control, or stop the task. These controls express
operator ownership of the session. They do not transform a hard block into an
approval and, while a security incident is active, ordinary resume is disabled
until the incident is resolved or the task is stopped.

\section{Narrow action approval}

\subsection{Information shown to the operator}

When policy returns \texttt{approval\_required}, execution pauses inside the
current tool step. The side panel docks an approval card above the composer so
the request remains adjacent to the live task. The card shows:

\begin{itemize}
  \item the proposed action name;
  \item a redacted current URL;
  \item the policy reason;
  \item the numeric risk score and named risk factors;
  \item a bounded, recursively redacted parameter preview; and
  \item a countdown to the server-owned expiry.
\end{itemize}

The available choices are \emph{Allow once}, \emph{Allow similar} when the
server explicitly offers a reusable scope, and \emph{Deny}. The API cannot
manufacture a reusable grant: a \texttt{similar} response is downgraded to
\texttt{once} if the original request did not permit it. This matters because
the interface is exposed code, whereas the original policy decision contains
the evidence needed to decide whether reuse is safe.

\subsection{Scope, expiry, and revalidation}

A one-time approval applies only to the pending proposal. A reusable grant is
still narrow and temporary. It binds the action type, target origin and
navigation path, capability set, taint class, sensitive-data categories,
and target-element fingerprint. Code, file writes, and uploads additionally
bind strict parameter or content evidence. A separate run-keyed parameter
digest joins audit records. The grant expires after a configured duration or
use count and is not transferable to another run.

After either a user response or grant match, \systemname\ repeats the complete
live preflight. It re-reads the current URL and target, reconstructs risk,
re-evaluates policy, and compares the new scope with the approved scope. A
navigation, DOM replacement, parameter or scope change, newly blocking or
handoff policy outcome, newly discovered sensitive value, unavailable security
context, or stop signal prevents execution. A relaxed policy may allow the
unchanged action; the earlier approval cannot authorise a broader one. The
terminal audit identifies
\texttt{approval\_revalidation} rather than implying that the user denied the
original action.

This ordering closes a time-of-check/time-of-use gap. A consent card describes
a specific state, but the browser remains mutable while the operator reads it.
Approval therefore authorises a bounded opportunity to try the action after
revalidation; it does not promise that the old action will execute.

\subsection{Prompt budget}

The default policy permits at most four approval prompts in a run. Once the
budget is exhausted, the runtime stops generating additional cards and follows
the configured fail-closed outcome, either blocking or transferring control.
The budget is a technical bound on interruption frequency. It should not be
misread as an empirically optimal number of prompts or evidence that four
prompts avoid fatigue.

\section{Direct human takeover}

\subsection{Why takeover is not approval}

Some tasks contain steps that should be completed by a person even if the
overall task remains legitimate. Examples include CAPTCHA completion and
entering a password, one-time code, access token, payment detail, private key,
or other protected value. If the agent were allowed to type that value after a
click on ``approve'', approval would merely authorise the very disclosure that
the guard identified. \systemname\ instead withholds the proposed agent action
and hands the live browser to the operator. This changes who performs the
entry; the later-observation limitation below prevents a stronger
non-interference claim.

Takeover can begin voluntarily when the agent requests help for a human step,
or mandatorily when policy detects a protected channel. The deterministic
guards cover indexed form fields, free-text keystrokes, relevant page-code
semantics, secret-looking upload paths, and readable local upload content
within a size bound. Missing, remote, binary, oversized, or unreadable uploads
are treated as unclassifiable and handed over rather than assumed safe. This
is specified channel coverage, not a claim of universal data-loss prevention.

\subsection{In-page ownership and input arbitration}

The content script renders a compact banner and page-edge signal from a closed
Shadow DOM. During ordinary autonomous work it states that the agent is
working and offers \emph{Take over}, \emph{Open chat}, and \emph{Stop}. A
transparent input shield blocks real keyboard input while the agent is
thinking, then drops during the agent's own action window so automation is not
intercepted. It is a cooperative cue that reduces overlap, not an isolation
boundary. Pause or takeover may wait for the in-flight BrowserUse step, and
Stop can prevent future governed execution but cannot undo a browser action
already executing.

During takeover, the visual state changes to ``Your turn'', the shield is
removed, and the banner provides a countdown, an optional note, and
\emph{Done -- continue}. The same request appears in the side panel. Either
surface can return control because both resolve the same opaque server request.
The note is operator-supplied agent context: its raw text returns to the model,
while events retain only a recursively redacted copy, so protected values must
not be placed in the note. The intercepted agent action is not executed, but
the prototype does not prove page-only secrecy after handback. BrowserUse
suppresses password-field values; non-password one-time-code, payment, and text
fields may reappear in a later DOM/accessibility observation.
For a sensitive-domain handover, successful completion clears only the matched
domain pattern for that run; it does not clear prompt-injection taint, override
hard policy, or authorise the intercepted disclosure.

When control returns, the agent observes the updated page state and replans.
Timeout and cancellation are explicit non-authorising outcomes; stop and
shutdown cause cancellation. If stop and handover completion race in the event
loop, stop wins before privileged execution.

\section{Unresolved-observation security incident}

\subsection{Quarantine before the model call}

The most security-sensitive intervention occurs when a detected current-page
observation cannot be sanitised with safe attribution or remains detector
positive after removal. Passing that content to the model with a warning would
still expose the decision-maker to the suspected attack. The current runtime
therefore withholds the observation before the provider call, records whether
the failure was \texttt{residual} or \texttt{unattributable}, and creates a
\texttt{security\_incident} control request.

The security-incident card is intentionally different from an approval card:

\begin{itemize}
  \item it contains no attacker text and accepts no operator note;
  \item it offers no way to approve or trust the hidden observation;
  \item it explains that no action was derived from the withheld observation;
  \item it remains visible across reconnects and even when optional everyday
  in-page visuals are disabled; and
  \item its safe choices are to inspect or change the page and request a fresh
  observation, or to stop the task.
\end{itemize}

The server and side panel retain a persistent incident state rather than an
ordinary ``working'' status. The in-page overlay mirrors it on a best-effort
basis but can be suppressed by hostile page code. The side-panel card uses an
assertive alert role, text labels in addition to colour, a countdown, and
explicit wording that re-checking does not approve the hidden content.
Reduced-motion preferences disable attention animations without removing the
side-panel status.

\subsection{Fresh-observation invariant}

Selecting \emph{Re-check page} resolves only the wait on the discarded
observation. BrowserUse begins a new step and captures a fresh representation
before any subsequent model call. If suspicious content remains, the new
observation is independently detected and can open another incident. If the
operator navigated away or removed it, processing continues from the new page.
The original prompt cannot be resumed or annotated into safety.

The same behaviour is not used for every short-lived wrapped value. An
unresolved page-extraction or history value is removed in full within a
bounded wrapper because opening an interactive browser incident for an old or
nested value can deadlock the model call or create an endless incident loop.
The distinction is based on recoverability: the operator can inspect and
change the current browser page, whereas an internal extracted or historical
string is safely omitted.

\section{Side panel, dashboard, and explanatory audit}

\subsection{Chat-first supervision}

The side panel is organised around one continuous task thread rather than a
separate security console. Persisted and live user text, agent step summaries
with proposed-action chips, non-allow policy decisions, detector alerts,
control requests and returns, and task outcomes appear in chronological
context. Pending security incidents are docked first, followed by ordinary
handovers and approvals, so the strongest unresolved state remains visually
dominant.

The header exposes run ownership controls and access to previous conversations,
in-page visuals, and the dashboard. The composer changes meaning with run
state: it starts or follows up on an idle conversation and queues steering for
the next thought while a task is active. An offline banner distinguishes a
local-service connection failure from an idle agent. Starting a new chat is
also a deliberate forced-reset path for a wedged run, which cancels any
outstanding authority rather than leaving orphaned requests.

Step summaries do not contain browser screenshot thumbnails. The ordinary
interactive model path is DOM/accessibility based, and the event fan-out
strips screenshot data before sending a step to the extension. Page titles and
URLs are redacted at event and trace boundaries. A locally retained screenshot
is a separate explicit audit option and is never the extension's live
observation channel. Likewise, displayed rationale-like summaries are exposed
model outputs and are not presented as private chain-of-thought.

\subsection{Activity and decision inspection}

The Activity dashboard lists completed runs with task, status, duration,
steps, model calls, and token totals. Expanding a run shows its result, a
decision log, and a within-run time partition. The UI labels policy checks plus
approval and incident waits as ``Enforcement added'', and labels the remainder
``Without enforcement''; these are accounting labels, not a matched
counterfactual. Handover wait is shown separately by implementation convention,
not because the same delay has been demonstrated without \systemname.
Enforcement changes later replanning, so the partition cannot estimate causal
overhead.

Decision chips count allowed actions, approvals, policy-triggered handoffs,
security pauses, blocks plus human denials,
approval/handoff timeouts, and total interventions. Non-zero chips deep-link
into the relevant trace events. Each non-allow policy decision in the decision
log can also receive a local yes/no \texttt{understood} annotation in a
researcher walkthrough. These ticks record whether the reviewer could explain
the displayed reason during inspection. They are not participant responses,
comprehension scores, or evidence of trust.

The Models view selects the primary and fallback providers, model identifiers,
sampling settings, and rate-limit pacing for the next task; API keys can be
written to the backend but are never read back into the client. The Policies
view exposes enforcement mode, domain and action rules, data-flow and
prompt-injection settings, trace retention, and run budgets. Live-evaluated
rule fields are re-read at the next model or action boundary; instantiated
model clients, timeouts, budgets, and schema pruning are fixed when a run
starts. A pending approval is revalidated against changed live rules, so an
obsolete consent cannot silently survive. These surfaces expose local control
without moving authority into the client.

\section{Audit and privacy invariants}

The interface is generated from the same events that support evaluation and
forensic review. Each model-proposed action entering the governed tool
boundary receives one \texttt{action\_audit:attempted} record before risky
preflight and one \texttt{action\_audit:completed} record for execution, tool
error, block, denial, timeout, cancellation, handover return, or exception. An
action identifier joins the pair. This prevents a preflight exception from
becoming an invisible governed attempt and prevents a card resolution from
being mistaken for execution.

The public record contains action class, verdict, policy rule, redacted reason,
capabilities, taint labels, risk factors, policy and preflight latency, target
origin, and a run-keyed parameter digest. Detection events contain score,
severity, signal names, and a finding digest, not the injection payload.
Sanitisation events contain branch and count telemetry, not removed text.
Event and history fields receive deterministic recursive pattern redaction
across nested parameters, URLs, free text, approval previews, operator notes,
errors, conversation history, and trace views. This is coverage for recognised
patterns, not proof that arbitrary sensitive or attacker-derived echoes are
absent. Step summaries remain exposed model output and may themselves be
sensitive.

Run-keyed HMAC digests correlate attempted and terminal records without
retaining the digested values or linking equal values across runs. Persistent
traces default to no screenshots and no free-form model rationale. Before
returning stored events, the Activity API again removes screenshots and legacy
untrusted page titles as defence in depth. The current runner also writes a
bounded prefix of steering text to its local process log; this path sits outside
trace redaction and prevents a categorical no-raw-free-text persistence claim.

\section{Failure semantics and technical verification}

\Cref{tab:interface-failures} summarises the principal pending-authority
failures. The invariant is that loss of required authority never becomes
implied consent: disconnect, expiry, stop, stale response, changed live
context, absent operator channel, and unresolved observation terminate or
remain pending without executing the withheld governed agent action. Direct
human and page-script effects remain outside that guarantee.

The verification strategy exercises these behaviours below the benchmark
level. Python tests cover request-manager lifecycle and pending-state
behaviour; scope restrictions; timeout and cancellation; selected
revalidation cases involving policy, DOM targets, uploads, and stop races;
guarded-channel takeover and upload classification; incident retry; recursive
redaction; and attempted/terminal audit pairing. Snapshot reconstruction is
implemented but lacks a full browser reconnect test. The TypeScript/Vite build
type-checks and packages the client; it does not prove Python--TypeScript
schema equivalence or test browser rendering, reconnect, accessibility, or
human behaviour.
\Cref{sec:implementation-verification} records the fresh submission-time gates.

These checks establish technical protocol properties, not human performance.
The service-backed extension presents selective-authority requests, persistent
incidents, and inspectable redacted evidence as a coherent control plane.
Whether operators notice the right information, decide well under time
pressure, or supervise several agents remains an empirical question; the
evaluation therefore uses contract tests as implementation evidence and makes
no human-outcome claim.
